This thesis presents an interpretable facial attractiveness analysis system that addresses the limitation of traditional holistic scoring approaches by decomposing beauty assessment into region-specific contributions. The system employs a dual-stream co-attention deep learning architecture combining MobileNetV2-based networks to process RGB facial images and 7-channel semantic parsing maps simultaneously. Using MediaPipe Face Mesh for precise detection of 478 facial landmarks, the system segments the face into eight anatomical regions (eyes, eyebrows, nose, mouth, skin, hair) and applies occlusion-based sensitivity analysis to quantify each feature's individual contribution to overall attractiveness. By systematically blurring specific facial regions and measuring the resulting score changes, the system reveals which features contribute positively or negatively to perceived beauty. Trained on the SCUT-FBP5500 dataset of 5,500 annotated facial images, the model achieves state-of-the-art performance with SRCC of 0.912 and PCC of 0.921 through 5-fold cross-validation. The system outputs both an overall attractiveness score (0-10 scale) and individual feature scores, accompanied by intuitive color-coded heatmaps that visualize contributions overlaid on the facial image. With a compact architecture of 7.9 million parameters and inference time of approximately 320 milliseconds per image, the system enables practical applications in personalized beauty recommendations, cosmetic consultation, aesthetic surgery planning, and virtual makeover systems, bridging the gap between automated beauty analysis and actionable aesthetic guidance.
